\documentclass[11pt]{article}
\usepackage[margin=0.9in]{geometry}
\usepackage{amsmath,booktabs,hyperref}
\title{Review of the Li-style ordinary-distance collision update}
\author{Pan Dynamics Collision Avoidance Research}
\date{October 2, 2026}
\begin{document}
\maketitle
\section{Consistency verdict}
The collision update in commit
\begin{center}\small\texttt{a073453f293ff7f47bcef4b9e92e096a9e13bbc9}\end{center}
follows the main sequence in Li et al., \emph{Real-Time Optimal Trajectory
Planning for Autonomous Driving with Collision Avoidance Using Convex
Optimization} (2023), DOI
\href{https://doi.org/10.1007/s42154-023-00222-7}{10.1007/s42154-023-00222-7}:
optimize the ordinary-distance dual, then fix its multipliers in the
trajectory problem. Sections 3.1--3.2 and equations (9)--(13) were read from
the user-provided local PDF; PDF page 6 was also visually inspected.
No author source code was examined.

For target inequalities $A(z-p_T)\le b$ and ego vertices $v_j(x)$, the new
implementation solves
\[
 \max_{\lambda\ge0,\alpha}\;\alpha-b^T\lambda,
 \qquad \alpha\le\lambda^TA(v_j(\bar x)-p_T)\quad\forall j,
 \qquad \|A^T\lambda\|_2\le1.
\]
It then fixes $\lambda^*$ in the four collision rows and linearizes the
rotated vertices with respect to ego pose. This explicitly implements a
finite-rectangle interpretation of the paper's set notation. The hypograph
variable, explicit vertex minimum, and yaw tangent are implementation details
beyond the printed formula; this review does not certify literal equivalence
to an unavailable author implementation.

The old signed-support replacement at overlapping poses, preferred direction
selection, and signed-direction MEX path have been removed. The ordinary
Euclidean set distance is nonnegative; overlap produces distance zero rather
than negative penetration. Importantly, the new multipliers are still fixed
within each trajectory solve. This update does not introduce joint online
optimization of the separating direction and trajectory.

\section{Important difference in the surrounding optimization}
Paper equation (13) includes a nonnegative collision slack $s_{k,m}$ in each
listed collision constraint, with a cost penalty. It does not describe this
project's hard completion tail. The project retains lexicographic PCBF/CLF
objectives, nonlinear Fiala-based affine dynamics, explicit finite-body yaw
linearization, and hard collision constraints after the primary horizon.
Thus adoption of the collision-dual sequence is supported; a claim that the
entire controller is the paper's algorithm would not be supported.

\section{Overlap degeneracy verified on the prior failing seed}
The saved flow-guided seed from planning time 0.85 s in the old
15-m/s turning-crossing run has overlapping rectangles at predicted time
1.65 s. The current ordinary-distance implementation evaluated at that same
pose returns
\[
 \|n\|_2=1.1261\times10^{-10},\qquad
 d_{\rm dual}=-7.42\times10^{-11}\ \mathrm m,
 \qquad \max_{j,k}|J_{jk}|=2.20\times10^{-10}.
\]
The tiny negative distance is numerical tolerance around the ordinary
zero distance. The conic solver returns flag 1. For the configured coordinate
correction box $|\Delta(x,y,\psi)|\le(2.5,2.5,0.25)$, even the independent
coordinate-wise upper bounds
\[
 g_j(\bar x)+|J_j|(2.5,2.5,0.25)^T
\]
are approximately $-0.10$ m for all four buffer-inclusive rows. This bound
ignores dynamics and other constraints, so the complete hard local problem
cannot repair this particular frozen near-zero-dual row.

For strictly overlapping convex bodies, the ordinary distance has the
zero-multiplier optimum; freezing zero multipliers removes the direction
that could separate them. A relaxed row reduces to $s\ge d_{\min}$ and
provides no state-dependent avoidance pressure. A hard positive-buffer row
reduces to $0\ge d_{\min}$, which is infeasible. Numerical multipliers in the
probe are near zero rather than identically zero, but the coordinate bound
establishes the same local obstruction. At contact the dual can be nonunique;
this review does not generalize the strict-overlap statement to every contact.

This finding concerns the compatibility of an overlapping search seed with
hard continuation constraints. It does not imply that ordinary distance duality
is mathematically incorrect, that the physical scenario is unavoidable, or
that the old signed-distance construction should be silently restored.

\section{Checks and experimental scope}
An independent rerun of \texttt{ordinaryDistanceDualTest} passed all 16 cases.
The supplied implementation record
\path{report/ORDINARY_DISTANCE_DUAL_20261002.json}
reports a combined 224 passing focused checks, after replacement of a former
restoration-success test by a hard-overlap failure test. Its independently
recorded 15-m/s \emph{crossing} scenario executes 17 of 19 requested holds,
then stops at 0.85 s with \texttt{pcbfNoNumericalResult}. Those results were
read, not rerun in this review, and concern crossing rather than the historical
turning-crossing seed used for the geometric probe. Test success therefore
does not establish scenario success.

Only source comparison, local-paper reading, focused geometry tests, and the
single-pose probe were performed. No complete controller campaign, runtime
qualification, production-source edit or configuration change was performed.
The companion directory records the MATLAB review script, numeric outputs,
test table and source hashes. The research conclusion is that the intended
distance-dual sequence is implemented, while its strict-overlap degeneracy
requires explicit treatment in the surrounding initialization/restoration
and hard-tail design before collision-free performance can be inferred.
\end{document}
